% Frequently asked questions about the KLXX manuscript.
% Standalone: compiles on its own and reads main.aux for cross-references,
% so figure, table, theorem and equation numbers track the manuscript.
% Link colors are inherited from preamble.tex, as in main.tex.
\documentclass[11pt]{article}
% Preamble for the X-regularization manuscript.
\usepackage[letterpaper, margin=1in]{geometry}
\usepackage{setspace}
\usepackage{amsmath,amssymb}
\usepackage{graphicx}
\usepackage{subcaption}
\usepackage[amsmath,thmmarks,framed]{ntheorem}
\usepackage{authblk}
\usepackage{tikz}
\usetikzlibrary{arrows.meta,calc,fit,positioning}
\usepackage{tcolorbox}
\usepackage{booktabs}
\usepackage{multirow}
\usepackage{placeins}
\usepackage{enumerate}
\usepackage{xcolor}
\definecolor{darkgreen}{rgb}{0.0,0.5,0.0}
\usepackage[colorlinks,linkcolor=red,citecolor=darkgreen,urlcolor=blue]{hyperref}
\hypersetup{
  pdftitle={Forward KL Can Be Better: Normalizing Flow Boltzmann Sampling with Log-Ratio Variation},
  pdfauthor={Qi Feng, Rongjie Lai, Di Qi, Xuda Ye}
}
\usepackage[round,authoryear]{natbib}

\usepackage[ruled,linesnumbered]{algorithm2e}
\SetKwInput{KwInput}{Input}
\SetKwInput{KwOutput}{Output}

\newtheorem{remark}{Remark}
\newtheorem{lemma}{Lemma}
\newtheorem{theorem}{Theorem}
\newenvironment{proof}{\par\medskip\noindent\textit{Proof.}\ }{\hfill$\square$\par\medskip}

\newcommand{\Le}{\leqslant}
\newcommand{\Ge}{\geqslant}
\newcommand{\D}{\mathrm{d}}
\newcommand{\KL}{\mathrm{KL}}
\newcommand{\X}{\mathrm{X}}
\newcommand{\ESS}{\mathrm{ESS}}

\theoremstyle{plain}

\usepackage{xr}
\externaldocument{main}

% preamble.tex sets linkcolor, citecolor and urlcolor but not filecolor or
% runcolor, whose hyperref defaults are cyan; cross-document references use
% those, so pin every key to the manuscript's scheme.
\hypersetup{
  pdftitle={Answers to frequently asked questions},
  colorlinks=true,
  linkcolor=red,
  citecolor=darkgreen,
  urlcolor=blue,
  filecolor=red,
  runcolor=red,
}

\title{Answers to frequently asked questions}
\author{}
\date{}

\newcounter{qa}
\newcommand{\QA}[2]{%
  \par\bigskip
  \noindent\textbf{Q\theqa.}~#1\par
  \medskip
  \noindent\textbf{A\theqa.}~#2\par
  \medskip
}
\newcommand{\Q}[1]{\refstepcounter{qa}\QA{#1}}

\begin{document}
\maketitle

\noindent
This note collects questions raised about
\emph{Mode Coverage in Normalizing Flow Boltzmann Generators via Log-Ratio
Variation} and the answers the manuscript supports.

\Q{How is mode coverage guaranteed? Is it by quench and temper, by the loss
with log-ratio variations, or by both?}
{Both, and neither alone suffices. Quench and temper supplies the candidate
modes: melting carries particles beyond the region the source covers, and
quenching drives them into the basins. What it does not supply is the
probability of each basin. It weights a well by the melted mass that falls into
it, so its outer wells carry far more than the target gives them, as
Figure~\ref{fig: 1d-qt} shows on the 1D Rastrigin potential. Feeding those
samples to a forward KL term as if they were target samples would therefore fit
the wrong distribution.

The log-ratio variation is the channel through which they can be used anyway. A
variation asks only that the log-ratio be flat under its weighting measure, so
the quench and temper samples enter as the points at which flatness is demanded,
never as an estimate of the mass they carry. The two are complementary: the
forward KL is accurate where its target samples land and blind elsewhere, while
the quench and temper distribution is inaccurate everywhere and blind nowhere.

That the variation alone is not enough either is visible in
Figure~\ref{fig: phi4-methods}: with the target surrogate as its only weighting
measure, $\KL$+$\X_\pi$ collapses onto one well of the tilted $\phi^4$ target at
every size and seed, returning $p_+\in\{0,1\}$. Coverage appears only when the
weighting measure carries the quench and temper samples.}

\Q{Does the propagation factor predict the error of a run?}
{No, and the manuscript does not use it that way. The factor summarizes weight
degeneracy along one accepted schedule. It is not a bound on the constant of
Theorem~\ref{thm: propagation}: the ordering of the inequality that produces it
runs the wrong way, and $\hat F_\Sigma$ is not a bound on $F_\Sigma$ at finite
$N$ either. A ratio of two factors is therefore not a prediction of a ratio of
two errors, and the one place where two factors are compared says so explicitly.

Nevertheless, $\hat F_\Sigma$ is computed directly from the effective sample
size of each stage, and it follows the summation form that the standard
non-asymptotic analysis of Feynman--Kac particle systems gives
\citep{delmoral2004feynman,chopin2020introduction}: a sum over the stage at
which the error enters, times a product over the stages it must still traverse.
The term injected at a stage is a second moment of the incremental weights, and
the stage effective sample size measures exactly that moment, so no estimation
step intervenes between the weight vectors a run already evaluates and the
reported number. The effective sample size is also the most widely used
diagnostic in importance sampling
\citep{liu2001monte,doucet2001introduction,martino2017rethinking}. A smaller
factor therefore reports a better-conditioned schedule, in the units the
analysis supplies, and nothing more.}

\Q{Why does the log-ratio variation use an $L^1$ difference rather than an
$L^2$ one, which is smoother?}
{Nothing in the analysis requires $L^1$. An $L^2$ difference is admissible, and
mode coverage would still be supplied by quench and temper, since what the
weighting measure contributes is the set of points at which flatness is
demanded and not the norm in which the deviation is measured.

The choice is practical. In early tests, not reported in the manuscript, the
$L^2$ form converged more slowly than the $L^1$ form, so we kept $L^1$. It also
carries the same scale as the log-ratio itself, so the coefficients and the
learning rate that suit the forward KL term suit the variation without
retuning; a squared difference changes that scale and would require the step
size to be adjusted with it. Finally, a mean absolute difference is the total
variation penalty familiar from image processing, where it is a long-established
choice \citep{rudin1992nonlinear}, rather than something introduced here.

The $L^2$ form is closely related to LDR-L2 \citep{schopmans2026ldr}, and
results for it can be found there.}

\Q{What is the difficulty in training without reference data?}
{The mass-covering objective needs samples from the target, and without data
they have to be manufactured during training. Every difficulty follows from
that. A sampler that manufactures them explores only where the current model
already places mass, so a mode the flow has not found generates no samples that
would teach it the mode exists, and the omission reinforces itself. The
manufactured samples are also biased, so the objective is a moving target rather
than a fixed one, which is what Theorem~\ref{thm: accuracy} addresses. The
alternative that needs no target samples, the reverse KL, is mode-seeking: it is
finite when the flow ignores part of the target, so it does not penalize the
omission at all. And manufacturing samples is not free, since each step must run
an annealing or MCMC sweep, and some methods must store or differentiate through
what it produces.

On molecular systems with a real force field the set of data-free works is
small. FAB \citep{midgley2022flow} learns alanine dipeptide from the
unnormalized density alone, and reports doing so before any other method.
Temperature-annealed Boltzmann generators \citep{schopmans2025temperature}
reach the physical target by lowering the temperature along a continuation.
Energy-only samplers exist that use no data at all, but have not been run on
force-field molecules: iDEM \citep{akhound2024iterated} is evaluated on Gaussian
mixtures, double wells and Lennard--Jones particle systems, the largest being
the $55$-particle system, and adjoint sampling \citep{havens2025adjoint} learns
a diffusion sampler from the energy by stochastic optimal control.

What separates these methods is not only how much they cost but when the cost
is paid. FAB maintains a prioritized replay buffer across training and
differentiates through its MCMC transitions; iDEM stores its reverse-SDE samples
in a replay buffer reused in the inner loop. A buffer lives inside the
optimization: it is written and read at every step, its entries must be
maintained and reweighted, and what it holds was produced by earlier and worse
models. The quench and temper set is instead an input to
Algorithm~\ref{alg: training}, built once before a stage is trained and only
subsampled thereafter, so it is paid offline and never interacts with the
gradient step. The generator keeps no buffer and differentiates through no MCMC
transition. Its additional cost over forward KL is that offline construction
together with the variation evaluations under the mixture measure, which on the
achiral molecules amounts to roughly twice the wall time.}

\Q{Why is there no comparison with other baseline methods, such as the reverse
KL or FAB?}
{This is a known gap in the manuscript. We have not run those comparisons, for
the reasons below, and we intend to add them in future work.

First, we did test the reverse KL and FAB on small 2D systems. The reverse KL
proved sensitive to singular potentials. Its integrand carries the term
$U(G^{-1}(y))$, so the gradient mixes the singularity of the target potential
with the Jacobian of the map, and training blew up readily; we abandoned it. For
FAB we found, in tests not reported in the manuscript, that at small batch sizes
its convergence was not better than that of the forward KL, FAB being generally
demanding of large batches. The forward KL is a simple and robust way to train a
flow, which is why we added the log-ratio variations to it rather than to
another loss.

Second, the methods differ structurally in ways that make a common standard hard
to define. Many require external data. FAB carries a replay buffer, at a cost
paid inside the optimization. iDEM does not track an explicit effective sample
size. TA-BG anneals along a different continuation coordinate, the temperature
rather than the regularization. Aligning these under one protocol is not a
matter of matching wall time.

Third, the manuscript does compare against forward KL and against
$\KL$+$\X_\pi$, and the latter is close to LDR-L1: at the same coefficient the
two agree within a factor of two, so running $\KL$+$\X_\pi$ covers that
comparison in substance. The forward KL is our baseline throughout, and most
comparisons are against these two.

Finally, we stress that KLXX is not a direct competitor to these methods. The
log-ratio variation is a principle for designing losses rather than one loss: it
is a flexible channel for bringing reference points into training, and it can be
added to any importance-weight-based method, FAB and the reverse KL included, by
forming the variation appropriate to each. What KLXX addresses is mode
discovery, by supplying candidate modes during Boltzmann generator training;
its accuracy relative to these methods remains to be explored. We intend this
paper to open the use of log-ratio variations rather than to settle it.}

\Q{Why are the molecular runs reported at a single seed?}
{The premise is narrower than it looks. The experiments that carry the accuracy
claims are multi-seed already. The $\phi^4$ tables report three seeds per
lattice size, against PT references themselves averaged over three independent
runs. The clock batch-size test averages four matched sampling seeds at each of
four batch sizes, and the particle-count test performs $2^{8-k}$ independent
rebuilds at each of seven sizes, equalized to $2.56\times10^6$ particles per
point. Single seeds appear in the molecular tables.

There the replication is across systems rather than across seeds, which is the
more informative axis. The alkane comparison spans ten matched settings over six
molecules, from methane to hexane, on two schedule types, and KLXX has the
smallest propagation factor in all ten. Those molecules differ in dimension, in
barrier structure and in the number of accepted stages: they are not one
experiment with the noise reseeded, and ten independent systems agreeing in the
same direction is not what chance produces. A per-setting error bar would add
little to a result already unanimous across the systems it is claimed for.

Multi-seed molecular runs are expensive, and the manuscript records in its
limitations where that cost falls. We would not describe a result that holds at
every size, on every molecule and in every matched setting as left in doubt for
want of a second seed on each.

A separate question is what a fixed compute budget buys, since KLXX costs
roughly twice the forward KL wall time. That comparison is untested and worth
running. It is not, however, a control on what the tables claim, which concerns
the conditioning of an accepted schedule at a given schedule length.}

\bibliographystyle{abbrvnat}
\bibliography{references}

\end{document}
